% Transcription — fonds Alexandre Grothendieck, shelfmark 98, batch 3
% (pages 41-43, the folder's last three pages). Established from the
% facsimile archives/batches/98/batch-03.pdf (a mirror of
% https://grothendieck.umontpellier.fr/98.pdf), per the skill
% transcribe-grothendieck, read from 300-600 dpi crops of that file (the
% tiles for this batch were still being cut when the pass was made). The
% batch file carries no cover sheet: PDF page k is folder page 40+k; the
% archivists' pencil numbers bottom-left (« 41 », « 42 », « 43 ») agree.
%
% Pass: Opus 5.5 (claude-opus-5-5), 2026-09-24 — first pass, unchecked against the pages by a human.
%
% The folder sits in the inventory group "69-89" « Géométrie et topologie
% combinatoire », whose subject does not match this folder's (as for folder
% 100); \dating{} carries the folder's own date verbatim, and nothing is
% concluded from the group.
%
% All three pages are transcribed; none is skipped. No page carries anyone
% else's text: all three are in his hand, so nothing is summarised.
%
% What the batch is.
%   41-42  Pencil and blue ink, fast, heavily reworked on 41. Headed
%          « Résultats … des » and « Avec les notations précédentes, mais
%          sans conditions de finitude »: a Lemma on base change for the
%          sheaves Ext. X flat over S, F and G S-flat, Ext^i(F,G) S-flat for
%          i <= N; then for every T over S there are canonical isomorphisms
%          Ext^i_{O_X}(F,G) (x)_S T -> Ext^i_{O_{X_T}}(F_(T), G_(T)). A long
%          diagonal addition in the left margin of 41 rewrites the hypotheses
%          (a local resolution L_* -> F|U of length <= N+1, the Coker of
%          Hom(L_N, G) -> Hom(L_{N+1}, G) S-flat); it is given as a
%          \marginal{}, its struck central block only as such. Proof: reduce
%          to S, X, T affine with rings A, B, A'; M, N the B-modules of F, G;
%          L_* a free complex, L_* (x)_A A' a resolution of M' in degrees
%          <= N; page 42: K^* = Hom(L_*, N), Ext^p_{B'}(M',N') = H^p(K^* (x)_A
%          A'), Ext^p_B(M,N) = H^p(K^*) for p <= N, flatness of the H^p(K^*)
%          and of the cokernel of K^N -> K^{N+1}, hence the base change
%          isomorphism; « cqfd ».
%   43     Black ink, fair: the plan of a chapter on « familles limitées »
%          of coherent sheaves and of complexes of coherent sheaves —
%          Prop. 1 to 7 and four corollaries, stated by their headings only
%          (finite unions, base change, H^p, relative H for proper f_i,
%          Ext^p, Hom, a relative Ext_{?,f}).
%
% Notation fixed for the batch. His underlined (sheaf) Ext, Hom and H are set
% \mathcal{E}xt, \mathcal{H}om, \mathcal{H}; the index of L and K, a small
% star or cross, is set * (L_*, K^*), as written. On 42 the index of Hom is
% read A where the argument would want B; flagged, not repaired. The sign in
% « Ext^i … x_S T » on 41 is a plain x on the page and is kept as \times.
%
% Cross-references. « les notations précédentes » (41) refers to pages before
% this batch (not identified here). No page numbers are cited.
\documentclass[11pt,a4paper]{article}
% Le préambule commun à toutes les transcriptions du fonds Grothendieck.
%
% Il tient en une idée : **l'incertitude se compose, elle ne se résout pas.**
% Une transcription qui lisse un mot douteux a détruit la seule chose pour
% laquelle on la faisait — un lecteur doit pouvoir savoir, à chaque endroit,
% ce qui était sur le feuillet et ce qui a été ajouté. Les six macros qui
% suivent sont donc l'appareil critique, et non de la décoration.
%
% Les mêmes commandes sont reconnues par `scripts/render.mjs`, qui produit la
% vue de lecture du site. Une macro ajoutée ici doit l'être aussi là-bas,
% faute de quoi le rendu échouera bruyamment — ce qui est voulu : un rendu
% silencieusement faux est pire qu'un rendu absent.

\ProvidesPackage{grothendieck}[2026/08/08 Transcriptions du fonds Grothendieck]

\usepackage{amsmath,amssymb,amsthm}
\usepackage{mathrsfs}
\usepackage{stmaryrd}
% Les diagrammes commutatifs sont partout dans ce fonds : c'est la langue
% même de la géométrie algébrique de Grothendieck, et une transcription qui
% les rendrait en prose ne transcrirait rien.
\usepackage{tikz-cd}
% Français par défaut — c'est la langue des feuillets — mais l'anglais est
% chargé aussi : la traduction et le résumé sont des documents anglais, et la
% typographie française (espaces avant les deux-points) y serait une faute.
% Ces documents basculent par \selectlanguage{english} en tête de corps.
\usepackage[english,french]{babel}
\usepackage{fontspec}
\usepackage{geometry}
\geometry{a4paper,margin=25mm,marginparwidth=28mm}
\usepackage{marginnote}
\usepackage{xcolor}
% ulem, not soul. `soul`'s \st and \ul are built on a character-by-character
% scan that XeLaTeX's font machinery breaks: the document fails with an
% undefined control sequence rather than degrading. ulem does the same two jobs
% and is Unicode-engine safe. `normalem` stops it hijacking \emph, which the
% transcriptions use for Grothendieck's own emphasis.
\usepackage[normalem]{ulem}
\usepackage{hyperref}
\hypersetup{colorlinks=true,linkcolor=black,urlcolor=black,pdfborder={0 0 0}}

% --- Métadonnées du lot -----------------------------------------------------
% Reprises telles quelles par le script de rendu, qui les affiche en tête de
% la vue de lecture. Elles fixent ce que le fichier prétend couvrir : sans
% elles, une transcription flotte sans point d'attache dans un fonds de seize
% mille feuillets.

\newcommand{\folder}[1]{\gdef\grothFolder{#1}}
\newcommand{\batch}[1]{\gdef\grothBatch{#1}}
\newcommand{\leaves}[2]{\gdef\grothFirst{#1}\gdef\grothLast{#2}}
\newcommand{\foldertitle}[1]{\gdef\grothFoldertitle{#1}}
\gdef\grothFolder{?}\gdef\grothBatch{?}\gdef\grothFirst{?}\gdef\grothLast{?}\gdef\grothFoldertitle{}

% --- Le résumé --------------------------------------------------------------
%
% Il ouvre la lecture modernisée, sous le titre « Résumé », et il est composé
% en petit. Ce n'est pas un ornement : c'est le seul passage du document qui
% ne soit pas des mathématiques, et un lecteur doit pouvoir le reconnaître
% avant de l'avoir lu — puis le sauter s'il connaît déjà le sujet, ou s'y
% arrêter s'il ne le connaît pas. Le filet qui le suit marque la reprise du
% corps.

\newenvironment{resume}
  {\small\setlength{\parskip}{0.45em}}
  {\par\medskip\hrule\medskip}

% --- Mots-clés ---------------------------------------------------------------
%
% En anglais, en clôture du résumé de la lecture modernisée : le vocabulaire
% moderne sous lequel chercher ce que les feuillets construisent. Le manifeste
% du site (`npm run manifest`) les extrait pour étiqueter la cote entière —
% cette ligne est la source unique des tags, il n'y a pas de fichier à côté.
\newcommand{\keywords}[1]{\par\smallskip\noindent
  {\footnotesize\textsc{Keywords} --- \textit{#1}}\par}

% --- L'appareil critique ----------------------------------------------------

% Le numéro de feuillet, en marge. C'est l'ancre qui permet de régler un
% désaccord sur une formule en regardant *un* feuillet plutôt qu'une chemise
% de six cents. Le site s'en sert aussi pour tourner les pages du fac-similé
% quand on fait défiler la transcription.
\newcommand{\leaf}[1]{%
  \marginnote{\footnotesize\color{gray!70}\textsf{#1}}%
  \label{leaf:#1}\ignorespaces}

% Illisible : on ne devine pas. Un mot inventé qui se lit comme les autres est
% le pire résultat possible.
\newcommand{\ill}{\textcolor{red!60!black}{[\,\ldots\,]}}

% Lecture douteuse : le mot est donné, et signalé comme incertain.
\newcommand{\uncertain}[1]{\uline{#1}}

% Ajout de l'éditeur — une lettre manquante, un mot sous-entendu.
\newcommand{\add}[1]{\textcolor{blue!50!black}{[#1]}}

% Biffé par Grothendieck lui-même. Ce qu'il a rayé fait partie du document :
% une rature dit souvent où la pensée a changé de direction.
\newcommand{\struck}[1]{\sout{#1}}

% Note du transcripteur, sur l'état matériel du feuillet.
\newcommand{\note}[1]{\par\smallskip{\small\color{gray!60!black}\textit{[#1]}}\par\smallskip}

% Note marginale de Grothendieck, distincte du corps du texte.
\newcommand{\marginal}[1]{\par\smallskip\noindent\hspace{1em}%
  {\small\color{gray!60!black}#1}\par\smallskip}

% --- Titre ------------------------------------------------------------------

\AtBeginDocument{%
  \begin{center}
    {\large\bfseries Fonds Alexandre Grothendieck --- cote n\textsuperscript{o}~\grothFolder}\\[2pt]
    {\itshape\grothFoldertitle}\\[4pt]
    {\small feuillets \grothFirst--\grothLast\ (lot \grothBatch)}\\[2pt]
    {\footnotesize\color{gray!60!black}Transcription établie d'après le fac-similé numérisé
     par l'Université de Montpellier.\\
     Les lectures incertaines sont soulignées, les illisibles notées [\,\ldots\,],
     les ajouts entre crochets.}
  \end{center}
  \vspace{1em}}

\endinput


\folder{98}
\batch{3}
\pages{41}{43}
\dating{[à partir de 1966-1967]}
\watermark{Édition de démonstration}
\foldertitle{Schémas de Hilbert. Normes : notes manuscrites (s.d.).}

\begin{document}

\section{[Changement de base pour les $\mathcal{E}xt$ de faisceaux plats]}
\note{titre de l'éditeur, entre crochets. La page 41 porte en tête, de sa
main, les mots donnés ci-dessous.}

\page{41}
Résultats \ill{} \ill{} des

\struck{Avec} Avec les notations précédentes, mais \add{sans} conditions de
finitude
\note{« sans » est écrit au-dessus de la ligne, « conditions de finitude »
en deux lignes sous la fin de la phrase.}

\textbf{Lemme.} Supposons que $X$ soit plat sur $S$, $F$ et $G$ plats sur
$S$, \struck{\ill{} par \ill{}} enfin $\mathcal{E}xt^i(F, G)$ plats sur $S$
\struck{$F$ \ill{} \ill{} \ill{} \ill{} \ill{}} pour $i \le N$. Alors pour
tout $T$ sur $S$, on a des isom.\ canoniques
\note{deux lignes biffées sous « $\mathcal{E}xt^i(F,G)$ plats sur $S$ » ;
on y devine $p \simeq x$ au-dessus d'un $F$ initial, rien d'autre. « Alors
pour tout $T$ sur $S$ » est encadré, et un trait vertical borde à gauche
tout l'énoncé des hypothèses.}
\[
  \mathcal{E}xt^i_{\mathcal{O}_X}(F, G) \times_S T \longrightarrow
  \mathcal{E}xt^i_{\mathcal{O}_{X_T}}(F_{(T)}, G_{(T)})
\]
\note{le signe entre $\mathcal{E}xt^i_{\mathcal{O}_X}(F,G)$ et $T$ est un
simple $\times$ indicé $S$.}

\marginal{On suppose que pour tout $x$ de $X$, il existe un voisinage $U$
de $x$ tel que, sur $U$, il y a une $F$-\struck{\ill{}} résolution
$L_* \to F|U$ de longueur $\le N+1$, avec $L_i$ \uncertain{loc.\ libres}
\ill{} \ill{} $\le N$, de type fini \ill{} \struck{\ill{} \uncertain{acycliques} \ill{}
$\le N+1$} \ill{} \ill{} \struck{\ill{} $\mathcal{H}^i =
\mathcal{H}om(L_*, M)$ \ill{} $\mathcal{H}om(L_N, N)$ \uncertain{sont}
\ill{} $\mathcal{H}om(L_{N+2}, N)$ \uncertain{soit}} et $S$-plats,
$\operatorname{Coker}\bigl(\mathcal{H}om(L_N, G) \to
\mathcal{H}om(L_{N+1}, G)\bigr)$ $S$-plat.}
\note{longue addition écrite en travers de la marge gauche, en oblique,
au crayon et à l'encre bleue ; aucun signe de renvoi, sinon le trait
vertical qui borde les hypothèses du lemme. Le bloc central, encadré
et barré de plusieurs traits, est donné comme biffé ; ses $M$ et $N$ sont
soulignés. Dans « longueur $\le N+1$ », $N$ récrit sur un mot noirci.
Les indices $L_N$, $L_{N+1}$ de la dernière ligne sont soulignés et peu
nets.}

\emph{Preuve.} Définition des isom., on \uncertain{suppose}.
$S$, $X$, $T$ sont affines, et on laisse au lecteur le soin de vérifier
que les isom.\ \uncertain{ainsi obtenus} « \uncertain{se recollent} ».

Soient $A$, $B$, $A'$ les \uncertain{anneaux} de $S$, $X$, $T$,
$M$, $N$ les $B$-modules qui définissent $F$, $G$, donc $B$ et $M$ sont
$A$-plats. \add{Soit $B' = B \otimes_A A'$.} Soit $L_*$ \ill{} $M$-
\struck{résolution} \uncertain{complexe} libre \struck{de type fini}
\ill{} de $B$-modules, \struck{$B$} \ill{} de degrés $\le N+1$,
\struck{\ill{}} \ill{} de type fini en \ill{} \ill{}
\struck{\ill{}} \ill{} les $L_i$, $M$ \ill{} \struck{$A'$} et acyclique en
dim.\ $\le N$ \struck{\ill{} \ill{} $A$ \ill{}} \ill{} les $A$-\ill{}
\struck{[\ill{}]}.
\note{ces lignes sont une suite de reprises superposées, au crayon et à
l'encre bleue ; seuls les mots et formules marqués ici se lisent, et leur
ordre n'est pas toujours sûr.}

Mais, \struck{$A$ (\ill{})} $L_* \otimes_A A'$ ($= L_* \otimes_B B'$)
\struck{est} en dim.\ $\le N$ \struck{\ill{}}, une résolution de
\[
  M \otimes_A A' = M \otimes_B B' = M',
\]
d'ailleurs formée de libres \struck{\ill{}} de type fini en \add{tt} dim.\
\ill{}
\note{le dernier mot de la page, noirci, ne se lit pas.}

\page{42}
Donc on a pour $p \le N$ \struck{\ill{}}
\[
  \operatorname{Ext}^p_{B'}(M', N') = H^p\bigl(\operatorname{Hom}(L'_*, N')\bigr)
  = H^p\bigl(\operatorname{Hom}_{\uncertain{A}}(L_*, N) \otimes_A A'\bigr)
\]
\note{le $N$ du dernier membre est récrit, peut-être sur $N'$ ; l'indice
de $\operatorname{Hom}$, ici et à la ligne suivante, se lit $A$ là où
l'on attendrait $B$ ; on le laisse tel quel.}

Soit donc $K^* = \operatorname{Hom}_{\uncertain{A}}(L_*, N)$, \struck{qui}
\ill{} \uncertain{longueur} \ill{} \ill{} et \ill{} $A$-\uncertain{plat}
\struck{[il \uncertain{suffira} \ill{} \ill{} $p < N$ \ill{} $N$ \ill{},
et pour $p = N$ \ill{} \ill{} \ill{}]}. Donc
\[
  \operatorname{Ext}^p_{B'}(M', N') = H^p(K^* \otimes_A A'), \qquad
  \operatorname{Ext}^p_{B}(M, N) = H^p(K^*) \qquad (p \le N).
\]
\note{« $p \le N$ » est écrit à droite des deux formules, séparé d'elles par
un trait vertical.}

D'ailleurs les $K^i$ sont $A$-\uncertain{libres}, les $H^p(K^*)$
\uncertain{aussi} pour $p \le N$, et \struck{\ill{}} le \uncertain{conoyau}
de $K^{N} \to K^{N+1}$ est $A$-plat. Il s'ensuit que l'homom.\ canonique
\[
  H^p(K^*) \otimes_A A' \longrightarrow H^p(K^* \otimes_A A')
\]
est un isom.\ en \add{tt} dimension, \struck{\ill{}}
en particulier pour $p \le N$ on trouve
\[
  \operatorname{Ext}^p_{B'}(M', N') \simeq \operatorname{Ext}^p_B(M, N) \otimes_A A'
\]
cqfd.
\note{l'exposant de $K^{N}$ est précédé d'une marque qui ne se lit pas ;
l'indice $A$ du dernier $\otimes$ est récrit sur une autre lettre.}

\section{[Familles limitées : plan]}
\note{titre de l'éditeur, entre crochets ; la page 43 n'en porte pas. Elle
donne une suite d'énoncés par leur seul intitulé.}

\page{43}
Famille limitée de faisceaux cohérents [ : \uncertain{depuis} \struck{\ill{}} bornés ]

Famille limitée de complexes de faisceaux cohérents (\ill{} \ill{})
\note{les deux premières lignes sont réunies par une accolade à gauche ;
le crochet « [ : … bornés ] » est relié par un trait à la parenthèse de la
seconde.}

\textbf{Prop.\ 1} Réunion finie de familles limitées est limitée.

\textbf{Prop.\ 2} Possibilité dans la définition de supposer
\struck{$S$ \ill{}} les \struck{familles} « familles
\uncertain{universelles} » plates sur $T$.

\textbf{Prop.\ 3} Changement de base \add{$S' \to S$} \struck{\ill{}}
transforme famille limitée en famille limitée. Réciproque si
$\operatorname{Im}(S' \to S)$ contient tous les points de $S$ intervenant
dans $F$.

\textbf{Prop.\ 4} $\mathcal{H}^p(K)$

\textbf{Prop.\ 5} $H^{s}_{f_1, \ldots, f_n}{}^{p}(K_1, \ldots, K_n)$, où
$f_i : X_i \to Z_i$ est propre, $K_i$ sur $X_i$.

\textbf{Cor.\ 1} \struck{\ill{}} Tors \uncertain{locaux} \struck{\ill{}}

\textbf{Cor.\ 2} Images inverses par $X' \to X$

\textbf{Cor.\ 3} Images directes et $R^p f_*$

\textbf{Prop.\ 6} $\mathcal{E}xt^p_{\mathcal{O}_X}(K, L)$ \quad [
\note{l'indice $\mathcal{O}_X$ est écrit entre cette ligne et la suivante ;
il peut aussi appartenir au $\mathcal{H}om$ du corollaire. Le crochet
ouvert reste sans suite.}

\textbf{Cor.} $\mathcal{H}om(F, G)$

\textbf{Prop.\ 7} $\mathcal{E}xt^p_{\uncertain{X, f}}(K, L)$
\uncertain{quand ça marche}

\textbf{Cor.} $\mathcal{E}xt^1_{\uncertain{X, f}}(F, G)$ en tout cas.
\note{l'indice de ces deux $\mathcal{E}xt$ se lit « e$x$, $f$ » ; la
première lettre peut être un $\mathcal{O}$ ou un $X$ cursif. Un trait
oblique, dans la marge gauche, part de « Prop.\ 6 » vers le haut.}

\end{document}
